\section*{第 \ref{ch:g4:geometry} 章 --- 直线与多边形}

\begin{solution}{exo:g4:geometry:1}
作图用到\cref{met:g4:geometry:draw}；小方框
表示 $d$ 与垂线之间的直角。
\end{solution}

\begin{solution}{exo:g4:geometry:2}
E：三条水平线段互相平行；每一条都与竖线垂直。
H：两条竖线段互相平行，
横线与两条竖线都垂直。N：两条竖线段互相平行；
斜线两者都不是。T：两条线段互相垂直。
Z：上下两条线段互相平行；
斜线两者都不是。
\end{solution}

\begin{solution}{exo:g4:geometry:3}
六边形；四边形；三角形；八边形；五边形。多边形的
顶点个数与边数相等：$6$、$4$、$3$、$8$、$5$。
\end{solution}

\begin{solution}{exo:g4:geometry:4}
五边形有 $5$ 条边。从 $A$ 出发，不相邻的顶点是 $C$ 和
$D$：两条对角线，$[AC]$ 与 $[AD]$。
\end{solution}

\begin{solution}{exo:g4:geometry:5}
$OA = 4$ cm（半径）；$OM = 4$ cm（圆上每一点到圆心的距离）；
$AB = 8$ cm（直径是半径的两倍）。
\end{solution}

\begin{solution}{exo:g4:geometry:6}
先沿一条边张开圆规，再把这个张口移到另外两条边上：
若圆规两脚的距离恰好与第二条边吻合，
则这两条边相等。
\end{solution}

\begin{solution}{exo:g4:geometry:7}
这条路径围成一个有 $5$ 条边的五边形 ---
但不是正五边形！
\end{solution}

\begin{solution}{exo:g4:geometry:8}
“垂直的两条直线一定相交”：正确 ---
它们恰好在直角处相交。

“平行的两条直线永不相交”：正确 --- 这正是平行的定义。

“一条直线可以既与某条直线平行，又与另一条直线垂直”：
正确 --- 在方格纸上，一条水平线与每一条水平线都平行，
并与每一条竖线都垂直。
\end{solution}

\begin{solution}{exo:g4:geometry:9}
直线 $d$ 与 $f$ 看上去平行 --- 而且它们确实总是平行的：
垂直于同一条直线的两条直线互相平行
（由\cref{prop:g6:lines:facts}正式确认这一点）。
\end{solution}

\begin{solution}{exo:g4:geometry:10}
一种正确的作图步骤是：“1. 画一个正方形 $ABCD$，边长为 $4$ cm。
2. 标出两条对角线的交点 $O$。3. 画圆，
圆心为 $O$、半径为 $2$ cm。”这个圆与每条边都相切，
切点恰为各边的中点。
\end{solution}

\begin{solution}{exo:g4:geometry:11}
四边形：$2$ 条对角线。五边形：$5$ 条。六边形：$9$ 条。每增加一个顶点，
对角线的条数增加得越来越多 --- 增长越来越快
（$+3$，然后 $+4$，\dots）。
\end{solution}
